\subsection{TOTALLY ORDERED SETS}

DEFINITION 9. Two elements \(x, y\) of a preordered set \(E\) are said to be comparable if the relation '' \(x \leqslant y\) or \(y \leqslant x\) '' is true. A set \(E\) is said to be totally ordered if it is ordered and if any two elements of \(E\) are comparable. The ordering on \(E\) is then said to be a total ordering, and the corresponding order relation a total order relation.

If \(x\) and \(y\) are elements of a totally ordered set, then \(x = y\) or \(x < y\) or \(x > y\) ; the negation of \(x \leqslant y\) is thus \(x > y\) .

An ordering on \(\mathbf{E}\) is a total ordering if and only if its graph \(\mathbf{G}\) satisfies the relation \(\mathbf{G} \cup \overline{\mathbf{G}} = \mathbf{E} \times \mathbf{E}\) , as well as the relations \(\mathbf{G} \circ \mathbf{G} = \mathbf{G}\) and \(\mathbf{G} \cap \overline{\mathbf{G}} = \Delta\) .

Examples

\begin{enumerate}
\def\labelenumi{(\arabic{enumi})}
\item
  Every subset of a totally ordered set is totally ordered by the induced ordering.
\item
  Let E be an arbitrary ordered set. The empty subset of E is totally ordered, and so is every subset of E consisting of a single element.
\end{enumerate}

* (3) The set \(\mathbf{R}\) of real numbers is totally ordered. *

\begin{enumerate}
\def\labelenumi{(\arabic{enumi})}
\setcounter{enumi}{3}
\tightlist
\item
  If A is a set which has at least two distinct elements, the set \(\mathfrak{P}(A)\) (ordered by inclusion) is not totally ordered, for if \(x \neq y\) , the subsets \(\{x\}\) and \(\{y\}\) are not comparable.
\end{enumerate}

A totally ordered set is also totally ordered with respect to the opposite ordering; it is a lattice and a fortiori both left and right directed.

PROPOSITION 11. Every strictly monotone mapping \(f\) of a totally ordered set \(\mathbf{E}\) into an ordered set \(\mathbf{F}\) is injective. If \(f\) is strictly increasing, \(f\) is an isomorphism of \(\mathbf{E}\) onto \(f(\mathbf{E})\) .

For \(x \neq y\) implies that either \(x < y\) or \(x > y\) ; hence

\[
f (x) <   f (y) \quad \text { or } \quad f (x) > f (y),
\]

so that \(f(x) \neq f(y)\) in either case. It remains to be shown that if \(f\) is strictly increasing, \(f(x) \leqslant f(y)\) implies \(x \leqslant y\) ; if not, we should have \(x > y\) , and therefore \(f(x) > f(y)\) .

PROPOSITION 12. Let E be a totally ordered set and let X be a subset of E. For an element b of E to be the least upper bound of X in E, it is necessary and sufficient that (1) b is an upper bound of X, (2) for every \(c \in E\) such that c \textless{} b, there exists \(x \in X\) such that \(c < x \leqslant b\) .

The second condition says that no element c \textless{} b is an upper bound of X, i.e., b is a minimal element of the set M of upper bounds of X; but this is the same as saying that b is the least element of M, since M is totally ordered (no. 10, Proposition 10).

